\documentclass[10pt]{article}

% ---------- Document Identity (update these only) ----------
\newcommand{\DocStage}{}
\newcommand{\DocTitle}{Your Road to Tier One SOF}
\newcommand{\ProgramName}{182-Week Civilian-to-Selection Readiness Pipeline}
\newcommand{\ProgramSubtitle}{}

% ---------- Page ----------
\usepackage[legalpaper,left=0.5in,right=0.5in,top=0.6in,bottom=0.45in]{geometry}

% ---------- Typography ----------
\usepackage[T1]{fontenc}
\usepackage[utf8]{inputenc}
\usepackage{libertinus}
\usepackage{microtype}
\usepackage{parskip}
\usepackage{ragged2e}
\usepackage{textcomp}
\AtBeginDocument{\color{black}}
\AtBeginDocument{\pagecolor{white}}
\AtBeginDocument{\justifying}

% ---------- Landscape pages ----------
\usepackage{pdflscape}

% ---------- Color ----------
\usepackage[table]{xcolor}
\definecolor{StageBlue}{HTML}{1F4E79}
\definecolor{StageBlueDark}{HTML}{163A5A}
\definecolor{LightGray}{HTML}{F2F4F7}
\definecolor{MidGray}{HTML}{D9DEE5}
\definecolor{RuleGray}{HTML}{C7CED8}
\definecolor{HeaderInk}{HTML}{000000}
\definecolor{Ink}{HTML}{000000}

% ---------- Utility ----------
\usepackage{enumitem}
\sloppy
\setlength{\emergencystretch}{2em}

% ---------- Boxes / UI ----------
\usepackage[most]{tcolorbox}

\tcbset{
  charterTitle/.style={
    enhanced,
    colback=StageBlue!4,
    colframe=StageBlue,
    boxrule=1.25pt,
    arc=3pt,
    left=16pt,right=16pt,top=14pt,bottom=14pt,
    borderline north={3.2pt}{0pt}{StageBlue},
    borderline west={4pt}{0pt}{StageBlue},
    borderline south={1.25pt}{0pt}{StageBlue},
    drop shadow={black!25!white},
  },
  charterRule/.style={
    enhanced,
    colback=LightGray,
    colframe=RuleGray,
    boxrule=0.85pt,
    arc=3pt,
    left=12pt,right=12pt,top=10pt,bottom=10pt,
    borderline west={3pt}{0pt}{StageBlue},
  },
  charterBadge/.style={
    enhanced,
    colback=StageBlue,
    coltext=white,
    colframe=StageBlueDark,
    boxrule=0.6pt,
    arc=2pt,
    left=6pt,right=6pt,top=3pt,bottom=3pt,
  }
}
\newtcbox{\Badge}{nobeforeafter,charterBadge}

% ---------- Lists ----------
\setlist[itemize]{leftmargin=1.5em, itemsep=0.25em, topsep=0.25em}
\setlist[enumerate]{leftmargin=1.8em, itemsep=0.25em, topsep=0.25em}

% ---------- TOC ----------
\usepackage{tocloft}
\setcounter{tocdepth}{3}
\setlength{\cftbeforesecskip}{3pt}
\setlength{\cftbeforesubsecskip}{2pt}
\setlength{\cftbeforesubsubsecskip}{0pt}
\renewcommand{\cftsecfont}{\bfseries}
\renewcommand{\cftsecpagefont}{\bfseries}
\renewcommand{\cftsubsecfont}{\normalfont}
\renewcommand{\cftsubsecpagefont}{\normalfont}
\renewcommand{\cftsubsubsecfont}{\normalfont}
\renewcommand{\cftsubsubsecpagefont}{\normalfont}
\setlength{\cftsecindent}{0pt}
\setlength{\cftsubsecindent}{1.25em}
\setlength{\cftsubsubsecindent}{2.8em}
\setlength{\cftsecnumwidth}{2.1em}
\setlength{\cftsubsecnumwidth}{2.8em}
\setlength{\cftsubsubsecnumwidth}{3.6em}

% Remove the built-in TOC heading so only the custom heading is shown
\makeatletter
\renewcommand{\tableofcontents}{\@starttoc{toc}}
\makeatother

% ---------- Header/Footer: page number only ----------
\usepackage{fancyhdr}
\pagestyle{fancy}
\fancyhf{}
\renewcommand{\headrulewidth}{0pt}
\renewcommand{\footrulewidth}{0pt}
\fancyfoot[C]{\thepage}

% ---------- Section formatting ----------
\usepackage{titlesec}

\titleformat{\section}
  {\Large\bfseries\color{Ink}}
  {\thesection}{0.6em}{}
  [\vspace{0.25em}\textcolor{StageBlue}{\rule{\linewidth}{0.8pt}}]

\titleformat{\subsection}
  {\large\bfseries\color{Ink}}
  {\thesubsection}{0.6em}{}

\titleformat{\subsubsection}
  {\normalsize\bfseries\color{Ink}}
  {\thesubsubsection}{0.6em}{}

\titlespacing*{\section}{0pt}{0.95em}{0.35em}
\titlespacing*{\subsection}{0pt}{0.65em}{0.25em}
\titlespacing*{\subsubsection}{0pt}{0.55em}{0.2em}

% ---------- Table engine ----------
\usepackage{tabularray}
\UseTblrLibrary{booktabs}

% ---------- tabularray: GLOBAL table treatment ----------
\SetTblrInner{
  cells = {fg=black, bg=white, valign=t},
}

% ---------- Header helpers ----------
\newcommand{\Hdr}[1]{\shortstack[c]{\strut #1\strut}}

% ---------- Lines ----------
\newcommand{\TitleLine}{\textcolor{StageBlue}{\rule{0.98\linewidth}{0.95pt}}}
\newcommand{\SubTitleLine}{\textcolor{RuleGray}{\rule{0.98\linewidth}{0.65pt}}}

% ---------- Continued on next page ----------
\newcommand{\ContinuedOnNextPageInline}{%
  \par
  \vfill
  \noindent\textcolor{RuleGray}{\rule{\linewidth}{1pt}}\vspace{-2pt}
  \noindent\hfill\textit{Continued on next page}\par\vspace{-10pt}
  \noindent\textcolor{RuleGray}{\rule{\linewidth}{1pt}}
  \clearpage
}

% ---------- PDF hyperlinks ----------
\usepackage[hidelinks]{hyperref}
\hypersetup{
  colorlinks=false,
  pdfborder={0 0 0},
  linktoc=all,
  pdftitle={\DocTitle},
  pdfauthor={\ProgramName}
}

% =========================================================
\begin{document}

\subsection*{Phase 10 Card Header}
\phantomsection
\addcontentsline{toc}{subsection}{Phase 10 Card Header}

\begin{itemize}
\item Phase \textbf{ID}: Phase 10
\item Phase Name: Phase 10 — Advanced Tactical Readiness Physical
\item Duration weeks: 12
\item Version: v1.0
\item Stage Alignment: Stage 5
\item Governing Status: Draft — Non-Deployable
\end{itemize}

\subsection*{1. Phase Mission}
\phantomsection

Phase 10 consolidates multi-domain readiness under strict validity and safety governance by producing gate-grade evidence only inside protected windows, while enforcing anti-stacking and reslot discipline so high-consequence tests remain comparable and admissible.

\subsection*{2. Phase Purpose}
\phantomsection

\begin{itemize}
\item Establish Phase 10 as the first phase where governance-gated underwater testing (\textbf{C15}) may be conditionally attempted, without weakening water safety doctrine.
\item Increase consequence and specificity of \textbf{RUN}, \textbf{RUCK}, \textbf{SWIM}, and \textbf{CAL} evidence while controlling fatigue contamination through protected windows and staged batteries.
\item Set up Phase 11 by proving repeatability and admissibility under controlled conditions rather than “one good day” outputs.
\end{itemize}

\subsection*{3. Entry Criteria}
\phantomsection

\begin{itemize}
\item Entry requires that the prior phase end gate was satisfied with admissible evidence (\textbf{VALID} sessions and \textbf{VALID} tests where required) and that no unresolved stop-rule events remain uncleared.
\item If evidence is incomplete or validity is compromised, default posture is \textbf{HOLD} and reslot to the next protected window per Stage 3 adjustment doctrine.
\item Underwater eligibility is not assumed at entry; it remains conditional and governance-gated within Phase 10.
\end{itemize}

\subsection*{4. Operating Constraints}
\phantomsection

\begin{itemize}
\item Cadence is fixed at 6 sessions per week across Phase 10. Frequency is not reduced to manage fatigue; dose and modality are downshifted instead.
\item Required session elements per Stage 1.F in correct order are mandatory:
  \begin{itemize}
  \item Safety self-check first
  \item Warm-up
  \item Mobility
  \item Main work
  \item Cardio
  \item Cooldown
  \item Recovery notes and required post-session notes per Stage 1.F
  \end{itemize}
\item Cardio minimum standard: minimum 30 minutes continuous per session.
\item 10,000 steps per day is a global requirement and must be logged.
\item Validity doctrine: admissible \textbf{VALID} evidence only for gates; \textbf{INVALID} tests provide no gate credit; missing required fields make evidence inadmissible.
\item Water governance is absolute; no unsupervised underwater or breath-hold exposure ever. Underwater testing is conditional and only permitted under Stage 2.C \textbf{C15} prerequisites; otherwise it is not scheduled and is inadmissible.
\end{itemize}

\subsection*{5. Primary Adaptations and Physiological Focus}
\phantomsection

\begin{itemize}
\item Improved locomotion readiness under controlled fatigue (\textbf{RUN} and \textbf{RUCK}) while preserving durability and foot integrity signals as gating constraints.
\item Sustained \textbf{CAL} repeatability under strict standards and artifact posture for admissible gate evidence.
\item \textbf{SWIM} unit-locked comparability and repeatability under controlled conditions.
\item Governance-gated underwater readiness as conditional capability, prioritized as safety-governed admissibility rather than a performance chase.
\item Recovery stability and evidence integrity (logging completeness, tool stability, environment stability) to support replicated gate decisions.
\end{itemize}

\subsection*{6. Primary Modalities}
\phantomsection

\begin{itemize}
\item Emphasized domains: \textbf{RUN}; \textbf{RUCK}; \textbf{CAL}; \textbf{SWIM}; \textbf{DURABILITY}; \textbf{RECOVERY}; \textbf{BODYCOMP} (supporting trend).
\item Conditional domain: \textbf{SWIM} underwater (\textbf{C15}) only when governance prerequisites are met; otherwise not scheduled.
\item De-emphasized/prohibited: maximal strength testing; new tests; impact escalation beyond governed progression; ruck load drift; any gate-grade testing outside protected windows.
\end{itemize}

\clearpage
\begin{landscape}

\subsection*{7. Weekly Structure}
\phantomsection

\begingroup
\scriptsize
\setlength{\tabcolsep}{2pt}
\renewcommand{\arraystretch}{1.12}

\begin{longtblr}{
  width=\linewidth,
  rowhead=1,
  colspec = {
    Q[l,wd=0.70in]
    X[l,1.90]
    X[l,1.25]
    X[l,3.15]
    X[l,2.90]
  },
  hlines,
  vlines,
  cells = {hsep=2pt, vsep=6pt, halign=l, valign=t},
  row{1} = {bg=StageBlue!15, fg=black, font=\bfseries, halign=l, valign=m},
  row{odd} = {bg=white},
  row{even} = {bg=LightGray},
}
\Hdr{Session \#} &
\Hdr{Primary Focus} &
\Hdr{Main Work Domains} &
\Hdr{Cardio Modality/Intent} &
\Hdr{Notes/Risk Controls} \\

1 &
Locomotion readiness, low-variance execution &
\textbf{RUN}; \textbf{DURABILITY} &
Modality: step-based walk or treadmill; Minutes: ≥30; RPE+Talk Test: RPE 3–5, full-to-short sentences; Continuity: continuous planned &
Preserve stable footwear/route/tool posture; downshift immediately on gait drift or hotspot escalation; no novelty near protected windows. \\

2 &
\textbf{CAL} quality and repeatability posture &
\textbf{CAL}; \textbf{DURABILITY} &
Modality: low-impact continuous (walk/bike/elliptical as needed); Minutes: ≥30; RPE+Talk Test: RPE 3–4, full sentences; Continuity: continuous planned &
\textbf{CAL} is technique-gated; artifact requirements govern gate admissibility; avoid fatigue stacking with major \textbf{RUN}/\textbf{RUCK}. \\

3 &
Load-tolerance support posture &
\textbf{RUCK}; \textbf{DURABILITY} &
Modality: low-impact continuous; Minutes: ≥30; RPE+Talk Test: RPE 3–5, full-to-short sentences; Continuity: continuous planned &
\textbf{RUCK} exposure is governed; do not introduce unlogged load changes; foot integrity is a gating constraint. \\

4 &
\textbf{SWIM} readiness and recovery control &
\textbf{SWIM}; \textbf{RECOVERY} &
Modality: low-impact continuous; Minutes: ≥30; RPE+Talk Test: RPE 2–4, full sentences; Continuity: continuous planned &
\textbf{SWIM} is unit-locked in testing; pool verification posture is mandatory for gate-grade attempts; if pool access unstable, reslot rather than improvise. \\

5 &
Recovery-biased consolidation &
\textbf{RECOVERY}; \textbf{DURABILITY} &
Modality: lowest-risk continuous modality; Minutes: ≥30; RPE+Talk Test: RPE 2–3, full sentences; Continuity: continuous planned &
Minimum Viable Session posture is the default downshift option; preserve cadence while lowering risk. \\

6 &
Pre-window stabilization or post-window re-stabilization &
\textbf{RECOVERY}; \textbf{DURABILITY}; \textbf{BODYCOMP} &
Modality: step-based or indoor substitute; Minutes: ≥30; RPE+Talk Test: RPE 2–4, full sentences; Continuity: continuous planned &
Treat as no-novelty enforcement day before protected windows; verify tool stability and environment unit identity readiness. \\

\end{longtblr}

\endgroup

\subsection*{8. Progression Rules}
\phantomsection

\begin{itemize}
\item One progression lever at a time rule: do not increase more than one of (a) weekly Cardio load, (b) longest Cardio-primary duration, (c) complexity/density of \textbf{CAL}/\textbf{RUN}/\textbf{RUCK} exposure in the same week.
\item Protect validity: no novelty in footwear, route class, pool unit/length assumptions, timing tools, or ruck configuration within the lead-in to protected windows.
\item Downshift posture: when \textbf{DURABILITY}/\textbf{RECOVERY} signals drift (pain, foot hotspots, sleep collapse, rising session RPE), downshift dose/modality rather than reducing weekly frequency.
\item No stacking make-ups: missed/invalid tests are reslotted; they are not recovered by stacking multiple major tests into one week.
\item Underwater: never progressed by “confidence.” It is either governance-authorized and executed per \textbf{C15} or it is not scheduled.
\end{itemize}

\subsection*{9. Deload and Test Placement}
\phantomsection

Stage 3 integrated calendar posture for Phase 10 is used:

\begin{itemize}
\item Week 1 baseline touchpoint exists.
\item Protected Deload+Test windows: Week 5 (mid-phase touchpoint), Week 10 (late-phase verification touchpoint), Week 12 (end-of-phase exit gate).
\item Week 11 is Transition (taper-aware stabilization; no new major test introductions).
\item Gate-grade attempts occur only in protected Deload+Test windows; build/intensify weeks do not host gate-grade batteries.
\end{itemize}

\end{landscape}

\clearpage
\begin{landscape}

\subsection*{10. Test Battery}
\phantomsection

\begingroup
\scriptsize
\setlength{\tabcolsep}{2pt}
\renewcommand{\arraystretch}{1.12}

\begin{longtblr}{
  width=\linewidth,
  rowhead=1,
  colspec = {
    X[l,1.55]
    X[l,3.25]
    Q[l,wd=0.95in]
    X[l,3.25]
    Q[l,wd=1.55in]
  },
  hlines,
  vlines,
  cells = {hsep=2pt, vsep=6pt, halign=l, valign=t},
  row{1} = {bg=StageBlue!15, fg=black, font=\bfseries, halign=l, valign=m},
  row{odd} = {bg=white},
  row{even} = {bg=LightGray},
}
\Hdr{Test Timing} &
\Hdr{Test \textbf{ID}/Name} &
\Hdr{Domain} &
\Hdr{Validity Conditions} &
\Hdr{Pass Threshold Stage 2 Tier} \\

Week 5 Deload+Test mid-phase touchpoint &
\textbf{C11} / \textbf{MET-2B-022} Swim 500 yd Time — No Fins or \textbf{C12} / \textbf{MET-2B-023} Swim 500 m Time — No Fins (unit locked) &
\textbf{SWIM} &
Pool length/unit verified; no fins; allowable strokes per protocol; warm-up logged; environment/tool fields complete; no stop-rule violations &
\textbf{INTERMEDIATE} \\

Week 5 Deload+Test mid-phase touchpoint &
\textbf{C3} / \textbf{MET-2B-010} Push-Ups 2-Minute Timed; \textbf{C4} / \textbf{MET-2B-011} Sit-Ups 2-Minute Timed; \textbf{C5} / \textbf{MET-2B-012} Pull-Ups Strict; \textbf{C6} / \textbf{MET-2B-013} Plank Forearm &
\textbf{CAL} &
\textbf{CAL} artifact evidence required for \textbf{VALID} gate use; warm-up logged; standards auditable; missing artifact = \textbf{INVALID} for gate use &
\textbf{INTERMEDIATE} \\

Week 10 Deload+Test late-phase verification &
\textbf{C14} / \textbf{MET-2B-019} 6-Mile Ruck Time Trial — 35 lb (dry load) &
\textbf{RUCK} &
Dry load verified; water logged separately; route or treadmill identity logged; footwear logged; warm-up logged; no stop-rule violations; missing required fields = \textbf{INVALID} &
\textbf{INTERMEDIATE} \\

Week 10 Deload+Test late-phase verification &
\textbf{C10} / \textbf{MET-2B-017} 5-Mile Run Time Trial or \textbf{C9} / \textbf{MET-2B-016} 3-Mile Run Time Trial (phase-specific selection) &
\textbf{RUN} &
Route measured or treadmill identity/settings logged; warm-up logged; no stop-rule violations; no non-safety pauses; missing required fields = \textbf{INVALID} &
\textbf{INTERMEDIATE} \\

Week 12 Deload+Test end-phase exit gate &
\textbf{C10} / \textbf{MET-2B-017} 5-Mile Run Time Trial (dominant \textbf{RUN} gate) &
\textbf{RUN} &
Same validity posture; staged to avoid fatigue contamination; reslot if invalid &
\textbf{INTERMEDIATE} \\

Week 12 Deload+Test end-phase exit gate &
\textbf{C14} / \textbf{MET-2B-019} 6-Mile Ruck Time Trial — 35 lb (dry load) (dominant \textbf{RUCK} gate) &
\textbf{RUCK} &
Same validity posture; staged to respect spacing intent; reslot if invalid &
\textbf{INTERMEDIATE} \\

Week 12 Deload+Test end-phase exit gate &
\textbf{C11} / \textbf{MET-2B-022} or \textbf{C12} / \textbf{MET-2B-023} Swim 500 (unit locked) &
\textbf{SWIM} &
Same validity posture; do not convert units; reslot if pool conditions compromise validity &
\textbf{INTERMEDIATE} \\

Week 12 Deload+Test end-phase exit gate &
\textbf{C3}–\textbf{C6} / \textbf{MET-2B-010}–013 \textbf{CAL} battery &
\textbf{CAL} &
Artifact posture mandatory for gate-validity; missing artifact = \textbf{INVALID} for gate use &
\textbf{INTERMEDIATE} \\

Week 10 or Week 12 conditional governed capability check &
\textbf{C15} / \textbf{MET-2B-025} Underwater Exposure Admissibility Flag; \textbf{C15} / \textbf{MET-2B-026} Underwater 25 m Time — One Breath &
\textbf{SWIM} &
Governance prerequisites mandatory: certified lifeguard + dedicated safety spotter; facility permission; no hyperventilation; abort discipline; missing governance = not scheduled and inadmissible &
\textbf{INTERMEDIATE} \\

\end{longtblr}

\endgroup

\end{landscape}
\clearpage

\subsection*{11. Exit Standards}
\phantomsection

\begin{itemize}
\item Objective gate to advance out of Phase 10 is produced in Week 12 Deload+Test window using admissible \textbf{VALID} evidence only.
\item Exit theme tier: \textbf{INTERMEDIATE}.
\item Minimum posture:
  \begin{itemize}
  \item All scheduled end-gate tests used for the phase exit decision are \textbf{VALID} (\textbf{INVALID} provides no gate credit and triggers reslot).
  \item Decision uses Stage 2 replication posture as governing rule; if replication is not met for a gate-critical metric, default posture is \textbf{HOLD} for that metric and reslot to regain valid replication evidence.
  \end{itemize}
\item Underwater:
  \begin{itemize}
  \item Underwater evidence may be used only if governance prerequisites are met and results are \textbf{VALID}; otherwise it is not scheduled and cannot be used for gate credit or “equivalency.”
  \end{itemize}
\end{itemize}

\subsection*{12. Risk Controls and Stop Rules}
\phantomsection

\begin{itemize}
\item Stage 0.5 and Stage 1.F control stop posture and escalation; this phase adds Phase 10-specific emphasis:
  \begin{itemize}
  \item Prevent stacked high-cost exposure in one week; stage batteries and prioritize decisive gate items; reslot overflow.
  \item Foot integrity is a controlling constraint for \textbf{RUCK} and \textbf{RUN} gate-grade attempts; hotspots/blisters trigger downshift and may force reslot.
  \item Environment hazards invalidate outdoor testing; substitute only when Stage 2 validity and Stage 2.E logging preserve comparability; otherwise reslot.
  \item Water governance remains absolute; underwater is conditional and supervised-only; if governance cannot be guaranteed, it is not scheduled.
  \end{itemize}
\end{itemize}

\clearpage
\begin{landscape}

\subsection*{13. Required Capabilities and Dependencies}
\phantomsection

Category-level only; feeds Stage 6 and Stage 7.

\begingroup
\scriptsize
\setlength{\tabcolsep}{2pt}
\renewcommand{\arraystretch}{1.12}

\begin{longtblr}{
  width=\linewidth,
  rowhead=1,
  colspec = {
    X[l,1.65]
    X[l,3.10]
    Q[l,wd=1.35in]
    Q[l,wd=1.20in]
    X[l,3.25]
  },
  hlines,
  vlines,
  cells = {hsep=2pt, vsep=6pt, halign=l, valign=t},
  row{1} = {bg=StageBlue!15, fg=black, font=\bfseries, halign=l, valign=m},
  row{odd} = {bg=white},
  row{even} = {bg=LightGray},
}
\Hdr{Dependency \textbf{ID}} &
\Hdr{Required Capability Category and Tier} &
\Hdr{Due-By week-in-phase} &
\Hdr{Owner} &
\Hdr{Fail Action} \\

\textbf{DEP-1C-007} \& \textbf{DEP-1C-008} &
7.11 Footwear Systems and Foot Care — Tier: must be stable and lockable for \textbf{RUN}/\textbf{RUCK} validity &
Week 1 &
Equipment Lead &
\textbf{HOLD} gate-grade \textbf{RUN}/\textbf{RUCK} attempts; substitute Cardio modality to protect tissue; reslot tests to next protected window once stable; document validity impact \\

\textbf{DEP-1C-017} &
7.07 Load-Bearing Rucking and Carry Systems — Tier: configuration stable and loggable &
Week 4 &
Equipment Lead &
Do not execute ruck time trials; reslot \textbf{RUCK} gate item; treat any attempt without dry-load verification as \textbf{INVALID} for gate use \\

\textbf{DEP-1C-006} &
7.18 Testing Timing Measurement and Course-Setup Tools — Tier: route measurement / treadmill identity / timing method stable &
Week 4 &
Coach Lead &
Do not execute gate-grade time trials; reslot; treat missing tool fields as \textbf{INVALID} for gate use \\

\textbf{DEP-1C-011} &
7.17 Aquatic Training and Water Safety Equipment — Tier: facility access and safety governance feasible &
Week 5 &
Coach Lead &
If pool verification or safety constraints fail, reslot \textbf{SWIM} tests; underwater is not scheduled without governance; no proxy credit \\

\textbf{DEP-1C-017} &
7.12 Recovery Mobility and Tissue-Care Implements — Tier: minimum recovery support available &
Week 1 &
Trainee &
Downshift to Minimum Viable Sessions when risk flags appear; if missing capability causes repeated drift, \textbf{HOLD} and stabilize before next gate window \\

\textbf{DEP-1C-015} &
7.09 Health Monitoring Diagnostics and Biometrics — Tier: baseline logging capability operational &
Week 1 &
Trainee &
Missing logs render evidence inadmissible; default \textbf{HOLD} and re-prove evidence chain before gate decisions \\

\end{longtblr}

\endgroup

\subsection*{14. Validity Requirements}
\phantomsection

\begin{itemize}
\item Evidence admissibility is governed by Stage 1.F and Stage 2: tests are \textbf{VALID}/\textbf{INVALID} only; invalid tests provide no gate credit.
\item \textbf{CAL} gate-validity requires compliant artifact evidence per Stage 2.C; missing/non-compliant artifact yields \textbf{INVALID} for gate use.
\item \textbf{RUN}/\textbf{RUCK} validity requires route/treadmill identity and settings (if treadmill), measurement integrity posture, warm-up logged, and no stop-rule violations.
\item \textbf{SWIM} validity requires unit-locked pool verification, no fins, allowable strokes, and rest governance compliance.
\item Underwater validity requires governance prerequisites; otherwise the test is not scheduled and is inadmissible.
\item Any compromised environment/tool change inside a protected window defaults to \textbf{INVALID} for gate use and triggers reslot per Stage 3.E.
\end{itemize}

\end{landscape}

\subsection*{15. Change Control Notes}
\phantomsection

\begin{itemize}
\item Allowed without patch: reslotting missed/invalid tests to the next protected deload/test window within the phase under Stage 3.E, without adding new tests or violating stacking intent.
\item Requires patch: changing gate test set, changing phase window placement, or redefining dependency posture.
\item No stacking make-ups: missed/invalid tests are not recovered by compressing multiple major tests into one week.
\end{itemize}

\subsection*{16. Compliance Checklist}
\phantomsection

\begin{itemize}
\item Phase 10 references only Stage 2 domain tags and Stage 2 C\#/\textbf{MET-2B} identifiers for gate tests.
\item Deload/test weeks are identified (Weeks 5/10/12) and Week 11 Transition posture is stated.
\item Weekly structure table includes six sessions and Cardio cell structure (Modality; Minutes; RPE + Talk Test; Continuity continuous planned).
\item Operating constraints include Cardio minimum 30 minutes continuous and 10,000 steps/day in body text.
\item Dependencies declared using Stage 7 categories and allowed owner taxonomy.
\item Governing status is correctly labeled Draft — Non-Deployable due to dependency \textbf{ID} placeholders affecting gating.
\end{itemize}

\end{document}